\clearpage
\begingroup
\let\clearpage\relax
\let\cleardoublepage\relax
\vspace*{\fill}
\chapter{虚拟内存}
\begin{center}
虚拟内存让每个进程都以为自己独占一整块连续地址空间：它把内存抽象成「按需装入的页」，
用磁盘充当后备存储，并以页面置换让有限的物理帧在多个进程间复用。
本章沿「按需分页 $\to$ 置换策略 $\to$ Belady 异常 $\to$ 抖动与工作集」的主线展开，
重点能用可计算的引用串手算缺页次数。
\end{center}
\vspace*{\fill}
\endgroup
\clearpage

\section{按需分页与缺页处理流程}

虚拟地址被划分为\textbf{页号}与\textbf{页内偏移}，由 MMU 借助页表把页号翻译成物理帧号。
按需分页的核心是：进程启动时并不把全部页调入内存，只有当指令或数据真正被访问、
且对应页不在内存时，才触发\textbf{缺页异常}把页调入。这种「懒惰装入」显著减小了驻留集与启动开销。

\subsection{地址翻译与页表项}

\begin{equation}
    \underbrace{\text{页号 } p}_{v-p \text{ 位}} \;\|\; \underbrace{\text{页内偏移 } d}_{p \text{ 位}}
    \xrightarrow{\ \text{页表}\ } \underbrace{\text{帧号 } f}_{\text{物理地址高位}} \;\|\; d
\end{equation}

页表中的每个页表项（PTE）除了帧号，还要携带若干状态位来支撑按需分页与置换：

\begin{table}[H]
\centering
\caption{页表项的典型字段}
\begin{tabulary}{\textwidth}{LLL}
\toprule
字段 & 含义 & 作用 \\
\midrule
有效位 (valid) & 该页是否在内存 & 为 0 时访问即触发缺页 \\
帧号 (frame) & 物理帧编号 & 地址翻译的结果 \\
脏位 (dirty) & 调入后是否被写过 & 决定换出时是否写回磁盘 \\
访问位 (reference) & 近期是否被访问 & 供 Clock/LRU 近似算法使用 \\
保护位 (protect) & 读/写/执行权限 & 越权访问产生保护异常 \\
\bottomrule
\end{tabulary}
\end{table}

\subsection{缺页处理流程}

触发缺页后，硬件先陷入内核，再由缺页处理程序完成调页。整个流程可概括为六步：
\textbf{trap} $\to$ \textbf{查页表} $\to$ \textbf{判定缺页} $\to$ \textbf{调页入内存} $\to$ \textbf{更新页表} $\to$ \textbf{重新执行指令}。

\begin{figure}[H]
\centering
\begin{tikzpicture}[font=\small, >=stealth, node distance=1.4cm and 2.4cm,
    every node/.style={draw, rounded corners, align=center, fill=blue!6, inner sep=5pt}]
    \node (start) {访问虚拟地址\\MMU 查页表};
    \node[below=of start] (valid) {有效位 $=1$ ?};
    \node[right=of valid, fill=green!12] (hit) {命中\\访问物理帧};
    \node[below=of valid, fill=red!12] (fault) {缺页异常 (trap)};
    \node[below=of fault] (find) {有空闲帧 ?};
    \node[right=of find, fill=orange!15] (load) {调入页面};
    \node[below=of find, fill=orange!15] (victim) {选牺牲页\\脏则写回};
    \node[below=of victim] (update) {更新 PTE / TLB};
    \node[below=of update, fill=green!12] (retry) {重新执行指令};

    \draw[->] (start) -- (valid);
    \draw[->] (valid) -- node[above]{是} (hit);
    \draw[->] (valid) -- node[right]{否} (fault);
    \draw[->] (fault) -- (find);
    \draw[->] (find) -- node[above]{是} (load);
    \draw[->] (find) -- node[right]{否} (victim);
    \draw[->] (load) |- (update);
    \draw[->] (victim) -- (update);
    \draw[->] (update) -- (retry);
    \draw[->, dashed] (retry.west) -- ++(-3.4,0) |- (valid.west);
\end{tikzpicture}
\caption{按需分页的缺页处理流程：查表 $\to$ trap $\to$ 调页 $\to$ 更新 $\to$ 重执行}
\end{figure}

缺页会带来磁盘 I/O，其代价远高于一次内存访问。设缺页率为 $p$，
则\textbf{有效访问时间}（EAT）为内存命中与缺页处理的加权平均：

\begin{equation}
    EAT = (1-p)\, t_{\text{mem}} + p\, t_{\text{fault}},
    \qquad
    p = \frac{F}{N}
\end{equation}

其中 $F$ 为缺页次数，$N$ 为总访问次数。$t_{\text{fault}}$ 通常在毫秒量级，
比 $t_{\text{mem}}$（约百纳秒）大四个数量级，因此\textbf{缺页率哪怕只有百分之几}
都会让性能急剧下降——这正是置换算法与工作集管理要解决的问题。

\section{页面置换算法：FIFO / LRU / 最优 OPT}

物理帧有限，当缺页且无空闲帧时必须\textbf{换出}某一页。选择牺牲页的规则即页面置换算法。
评价指标是给定引用串下的\textbf{缺页次数}。本节以经典引用串 $20$ 次访问为例：

\begin{center}
\texttt{7, 0, 1, 2, 0, 3, 0, 4, 2, 3, 0, 3, 2, 1, 2, 0, 1, 7, 0, 1}
\end{center}

\subsection{三种算法的直觉}

\begin{itemize}
    \item \textbf{FIFO}：淘汰\textbf{最早进入}内存的页，用队列实现；简单，但可能换出仍频繁使用的页，且存在 Belady 异常。
    \item \textbf{LRU}：淘汰\textbf{最久未被使用}的页，符合局部性；需要记录访问时序，硬件完全实现代价高，常用 Clock 近似。
    \item \textbf{OPT}（最优）：淘汰\textbf{未来最长时间不会被访问}的页；缺页次数理论下界，但需预知未来，仅作基准。
\end{itemize}

\subsection{帧数 $=3$ 的逐步模拟}

下表给出三种算法在相同引用串、相同 $3$ 个物理帧下的逐步状态（``—'' 表示空帧）。

\begin{table}[H]
\centering
\caption{FIFO（先进先出），共缺页 $15$ 次}
\begin{tabular}{rrrrrl}
\toprule
步 & 访问页 & 帧1 & 帧2 & 帧3 & 结果 \\
\midrule
 1 & 7 & 7 & — & — & 缺页 \\
 2 & 0 & 7 & 0 & — & 缺页 \\
 3 & 1 & 7 & 0 & 1 & 缺页 \\
 4 & 2 & 2 & 0 & 1 & 缺页（换出 7） \\
 5 & 0 & 2 & 0 & 1 & 命中 \\
 6 & 3 & 2 & 3 & 1 & 缺页（换出 0） \\
 7 & 0 & 2 & 3 & 0 & 缺页（换出 1） \\
 8 & 4 & 4 & 3 & 0 & 缺页（换出 2） \\
 9 & 2 & 4 & 2 & 0 & 缺页（换出 3） \\
10 & 3 & 4 & 2 & 3 & 缺页（换出 0） \\
11 & 0 & 0 & 2 & 3 & 缺页（换出 4） \\
12 & 3 & 0 & 2 & 3 & 命中 \\
13 & 2 & 0 & 2 & 3 & 命中 \\
14 & 1 & 0 & 1 & 3 & 缺页（换出 2） \\
15 & 2 & 0 & 1 & 2 & 缺页（换出 3） \\
16 & 0 & 0 & 1 & 2 & 命中 \\
17 & 1 & 0 & 1 & 2 & 命中 \\
18 & 7 & 7 & 1 & 2 & 缺页（换出 0） \\
19 & 0 & 7 & 0 & 2 & 缺页（换出 1） \\
20 & 1 & 7 & 0 & 1 & 缺页（换出 2） \\
\bottomrule
\end{tabular}
\end{table}

\begin{table}[H]
\centering
\caption{LRU（最近最久未使用），共缺页 $12$ 次}
\begin{tabular}{rrrrrl}
\toprule
步 & 访问页 & 帧1 & 帧2 & 帧3 & 结果 \\
\midrule
 1 & 7 & 7 & — & — & 缺页 \\
 2 & 0 & 7 & 0 & — & 缺页 \\
 3 & 1 & 7 & 0 & 1 & 缺页 \\
 4 & 2 & 2 & 0 & 1 & 缺页（换出 7） \\
 5 & 0 & 2 & 0 & 1 & 命中 \\
 6 & 3 & 2 & 0 & 3 & 缺页（换出 1） \\
 7 & 0 & 2 & 0 & 3 & 命中 \\
 8 & 4 & 4 & 0 & 3 & 缺页（换出 2） \\
 9 & 2 & 4 & 0 & 2 & 缺页（换出 3） \\
10 & 3 & 4 & 3 & 2 & 缺页（换出 0） \\
11 & 0 & 0 & 3 & 2 & 缺页（换出 4） \\
12 & 3 & 0 & 3 & 2 & 命中 \\
13 & 2 & 0 & 3 & 2 & 命中 \\
14 & 1 & 1 & 3 & 2 & 缺页（换出 0） \\
15 & 2 & 1 & 3 & 2 & 命中 \\
16 & 0 & 1 & 0 & 2 & 缺页（换出 3） \\
17 & 1 & 1 & 0 & 2 & 命中 \\
18 & 7 & 1 & 0 & 7 & 缺页（换出 2） \\
19 & 0 & 1 & 0 & 7 & 命中 \\
20 & 1 & 1 & 0 & 7 & 命中 \\
\bottomrule
\end{tabular}
\end{table}

\begin{table}[H]
\centering
\caption{最优置换 OPT，共缺页 $9$ 次（理论下界）}
\begin{tabular}{rrrrrl}
\toprule
步 & 访问页 & 帧1 & 帧2 & 帧3 & 结果 \\
\midrule
 1 & 7 & 7 & — & — & 缺页 \\
 2 & 0 & 7 & 0 & — & 缺页 \\
 3 & 1 & 7 & 0 & 1 & 缺页 \\
 4 & 2 & 2 & 0 & 1 & 缺页（换出 7） \\
 5 & 0 & 2 & 0 & 1 & 命中 \\
 6 & 3 & 2 & 0 & 3 & 缺页（换出 1） \\
 7 & 0 & 2 & 0 & 3 & 命中 \\
 8 & 4 & 2 & 4 & 3 & 缺页（换出 0） \\
 9 & 2 & 2 & 4 & 3 & 命中 \\
10 & 3 & 2 & 4 & 3 & 命中 \\
11 & 0 & 2 & 0 & 3 & 缺页（换出 4） \\
12 & 3 & 2 & 0 & 3 & 命中 \\
13 & 2 & 2 & 0 & 3 & 命中 \\
14 & 1 & 2 & 0 & 1 & 缺页（换出 3） \\
15 & 2 & 2 & 0 & 1 & 命中 \\
16 & 0 & 2 & 0 & 1 & 命中 \\
17 & 1 & 2 & 0 & 1 & 命中 \\
18 & 7 & 7 & 0 & 1 & 缺页（换出 2） \\
19 & 0 & 7 & 0 & 1 & 命中 \\
20 & 1 & 7 & 0 & 1 & 命中 \\
\bottomrule
\end{tabular}
\end{table}

\begin{table}[H]
\centering
\caption{三种算法缺页次数对比（引用串 20 次，帧数 3）}
\begin{tabulary}{\textwidth}{LLLL}
\toprule
算法 & 缺页次数 & 缺页率 & 备注 \\
\midrule
FIFO & 15 & 75.0\% & 实现最简单，可能产生异常 \\
LRU  & 12 & 60.0\% & 贴近局部性，硬件实现昂贵 \\
OPT  & 9  & 45.0\% & 理论最优，不可在线实现 \\
\bottomrule
\end{tabulary}
\end{table}

\section{Belady 异常}

直觉上增加物理帧应当减少缺页。但对 \textbf{FIFO}，
存在引用串使\textbf{帧数增多、缺页反而增多}，这一反直觉现象称为 \textbf{Belady 异常}。
经典反例引用串为：

\begin{center}
\texttt{1, 2, 3, 4, 1, 2, 5, 1, 2, 3, 4, 5}
\end{center}

\begin{table}[H]
\centering
\caption{FIFO 在 $3$ 个帧下缺页 $9$ 次}
\begin{tabular}{rrrl}
\toprule
步 & 访问页 & 帧 & 结果 \\
\midrule
 1 & 1 & 1 & 缺页 \\
 2 & 2 & 1\ 2 & 缺页 \\
 3 & 3 & 1\ 2\ 3 & 缺页 \\
 4 & 4 & 4\ 2\ 3 & 缺页（换出 1） \\
 5 & 1 & 4\ 1\ 3 & 缺页（换出 2） \\
 6 & 2 & 4\ 1\ 2 & 缺页（换出 3） \\
 7 & 5 & 5\ 1\ 2 & 缺页（换出 4） \\
 8 & 1 & 5\ 1\ 2 & 命中 \\
 9 & 2 & 5\ 1\ 2 & 命中 \\
10 & 3 & 5\ 3\ 2 & 缺页（换出 1） \\
11 & 4 & 5\ 3\ 4 & 缺页（换出 2） \\
12 & 5 & 5\ 3\ 4 & 命中 \\
\bottomrule
\end{tabular}
\end{table}

\begin{table}[H]
\centering
\caption{FIFO 在 $4$ 个帧下缺页 $10$ 次（比 $3$ 帧更多，即 Belady 异常）}
\begin{tabular}{rrrrrl}
\toprule
步 & 访问页 & 帧1 & 帧2 & 帧3 & 帧4 \\
\midrule
 1 & 1 & 1 & — & — & — \\
 2 & 2 & 1 & 2 & — & — \\
 3 & 3 & 1 & 2 & 3 & — \\
 4 & 4 & 1 & 2 & 3 & 4 \\
 5 & 1 & 1 & 2 & 3 & 4 \\
 6 & 2 & 1 & 2 & 3 & 4 \\
 7 & 5 & 5 & 2 & 3 & 4 \\
 8 & 1 & 5 & 1 & 3 & 4 \\
 9 & 2 & 5 & 1 & 2 & 4 \\
10 & 3 & 5 & 1 & 2 & 3 \\
11 & 4 & 4 & 1 & 2 & 3 \\
12 & 5 & 4 & 5 & 2 & 3 \\
\bottomrule
\end{tabular}
\end{table}

结论：FIFO 会发生 Belady 异常，因为它\textbf{不具备栈性质}。
若把 $n$ 个帧的驻留集合记为 $S_n(t)$，栈性质要求
$S_n(t)\subseteq S_{n+1}(t)$（任何时刻小内存的页集合是大内存页集合的子集）。
\textbf{LRU 与 OPT 都是栈式算法}，因此\textbf{永不发生 Belady 异常}，缺页次数随帧数单调不增。

\begin{equation}
    \text{栈性质：}\ S_n(t)\subseteq S_{n+1}(t)\ \Longrightarrow\ F_n(t)\ge F_{n+1}(t)
\end{equation}

\section{抖动与工作集模型}

\subsection{抖动（Thrashing）的成因}

当多道程序度过高、进程的\textbf{工作集之和超过可用物理帧总数}时，会发生连锁反应：
进程频繁缺页 $\to$ 内核忙于换入换出 $\to$ 就绪队列中进程长期得不到完整驻留 $\to$
缺页更频繁。此时 CPU 大部分时间在等待磁盘，\textbf{实际有效利用率反而骤降}，
这种现象称为\textbf{抖动}。

\begin{figure}[H]
\centering
\begin{tikzpicture}
\begin{axis}[
    width=0.82\textwidth, height=6.2cm,
    xlabel={多道程序度（驻留进程数）},
    ylabel={CPU 利用率 (\%)},
    xmin=0, xmax=10, ymin=0, ymax=100,
    grid=both, grid style={gray!20},
    every axis plot/.append style={thick}
]
\addplot[smooth, mark=*, blue] coordinates {
    (0,0) (1,20) (2,40) (3,60) (4,75) (5,86) (6,88) (7,80) (8,58) (9,32) (10,14)
};
\addplot[dashed, red] coordinates {(6,0) (6,100)};
\node[anchor=south west, red, font=\scriptsize] at (axis cs:6.05,10) {拐点};
\end{axis}
\end{tikzpicture}
\caption{CPU 利用率随多道程序度先升后降，下降段即抖动区}
\end{figure}

\subsection{工作集模型}

工作集模型用「最近一段时间内被访问的页」来刻画进程的\textbf{局部性}。
给定窗口大小 $\Delta$（以最近 $\Delta$ 次内存访问计），定义时刻 $t$ 的工作集：

\begin{equation}
    W(t,\Delta) = \{\, \text{最近的 } \Delta \text{ 次访问中出现过的页} \,\}
\end{equation}

工作集大小 $|W(t,\Delta)|$ 是进程在窗口内真正需要的帧数。系统的\textbf{局部性}体现在：

\begin{itemize}
    \item \textbf{时间局部性}：刚被访问的页很可能马上又被访问，故 LRU 有效。
    \item \textbf{空间局部性}：相邻地址的页常被连续访问，故预取与页大小选择有意义。
\end{itemize}

据此可判断抖动并给出分配准则：

\begin{equation}
    \sum_{i} |W_i(t,\Delta)| > M \;\Longrightarrow\; \text{抖动}
    \qquad
    \text{分配准则：帧}_i \ge |W_i(t,\Delta)|
\end{equation}

其中 $M$ 为全部物理帧数。若各进程工作集之和超过 $M$，则必须\textbf{挂起或换出（减少多道程序度）}
部分进程，直到剩余进程都能获得不小于其工作集的帧数，抖动才能消退。
窗口 $\Delta$ 太小会低估需求、太大又会把不相关的历史页计入，实践中常配合缺页率（PFF）动态调整。

\begin{equation}
    \text{缺页率 } p = \frac{F}{N}:
    \quad p > p_{\text{upper}} \Rightarrow \text{加帧},\qquad
    p < p_{\text{lower}} \Rightarrow \text{减帧}
\end{equation}
